% ****** Start of file apssamp.tex ******
%
%   This file is part of the APS files in the REVTeX 4.2 distribution.
%   Version 4.2a of REVTeX, December 2014
%
%   Copyright (c) 2014 The American Physical Society.
%
%   See the REVTeX 4 README file for restrictions and more information.
%
% TeX'ing this file requires that you have AMS-LaTeX 2.0 installed
% as well as the rest of the prerequisites for REVTeX 4.2
%
% See the REVTeX 4 README file
% It also requires running BibTeX. The commands are as follows:
%
%  1)  latex apssamp.tex
%  2)  bibtex apssamp
%  3)  latex apssamp.tex
%  4)  latex apssamp.tex
%
\documentclass[%
 reprint,
%superscriptaddress,
%groupedaddress,
%unsortedaddress,
%runinaddress,
%frontmatterverbose, 
%preprint,
%preprintnumbers,
%nofootinbib,
%nobibnotes,
%bibnotes,
 amsmath,amssymb,
 aps,
%pra,
%prb,
%rmp,
%prstab,
%prstper,
%floatfix,
]{revtex4-2}

\usepackage{graphicx}% Include figure files
\usepackage{dcolumn}% Align table columns on decimal point
\usepackage{bm}% bold math
%\usepackage{hyperref}% add hypertext capabilities
%\usepackage[mathlines]{lineno}% Enable numbering of text and display math
%\linenumbers\relax % Commence numbering lines

%\usepackage[showframe,%Uncomment any one of the following lines to test 
%%scale=0.7, marginratio={1:1, 2:3}, ignoreall,% default settings
%%text={7in,10in},centering,
%%margin=1.5in,
%%total={6.5in,8.75in}, top=1.2in, left=0.9in, includefoot,
%%height=10in,a5paper,hmargin={3cm,0.8in},
%]{geometry}

\begin{document}

\preprint{APS/123-QED}

\title{Spectral functions via novel cumulant expansions}% Force line breaks with \\
% \thanks{A footnote to the article title}%

% \author{Ann Author}
%  \altaffiliation[Also at ]{Physics Department, XYZ University.}%Lines break automatically or can be forced with \\
% \author{Second Author}%
%  \email{Second.Author@institution.edu}
% \affiliation{%
%  Authors' institution and/or address\\
%  This line break forced with \textbackslash\textbackslash
% }%

% \collaboration{MUSO Collaboration}%\noaffiliation

% \author{Charlie Author}
%  \homepage{http://www.Second.institution.edu/~Charlie.Author}
% \affiliation{
%  Second institution and/or address\\
%  This line break forced% with \\
% }%
% \affiliation{
%  Third institution, the second for Charlie Author
% }%
% \author{Delta Author}
% \affiliation{%
%  Authors' institution and/or address\\
%  This line break forced with \textbackslash\textbackslash
% }%

% \collaboration{CLEO Collaboration}%\noaffiliation

\date{\today}% It is always \today, today,
%              %  but any date may be explicitly specified
% \usepackage{authblk}
% \author[1]{Patryk Kozlowski}
% \affil[1]{Department of Chemistry and Chemical Biology, Harvard University, Cambridge, MA 02138, USA}


% \title{Theoretical Spectroscopy via Many-Body Perturbation Theory}
% \begin{abstract}
% An article usually includes an abstract, a concise summary of the work
% covered at length in the main body of the article. 
% \begin{description}
% \item[Usage]
% Secondary publications and information retrieval purposes.
% \item[Structure]
% You may use the \texttt{description} environment to structure your abstract;
% use the optional argument of the \verb+\item+ command to give the category of each item. 
% \end{description}
% \end{abstract}

%\keywords{Suggested keywords}%Use showkeys class option if keyword
                              %display desired
\maketitle

%\tableofcontents
\section{Introduction}
\label{sec:intro}
Electron-boson interactions are ubiquitous in condensed matter physics. For example, in core spectroscopy, the injection of a localized charge carrier couples to the plasmonic modes of the system, forming a fermionic polaron. This manifests in the observed spectrum as a sharp quasiparticle peak accompanied by satellite structures separated by the plasma frequency. In order to facilitate the interpretation of experiments, one may compute the spectral function, which is the direct observable in spectroscopy experiments.
% , 

Many-body perturbation theory (MBPT) is commonly used to obtain the spectral function, which is proportional to the imaginary part of the retarded single-particle Green's function $G$. MBPT obtains the Green's function through perturbative expansion in the bosonic propagator; in electronic structure, it is common to choose bosonic propagator to be the screened Coulomb interaction $W$. Doing so has the function of improving upon the mean-field answer at reasonable computational cost. The lowest-order term in this perturbative expansion is $G_0W_0$, which has been empirically successful at describing quasiparticles in systems with weak to intermediate coupling. However, $G_0W_0$ famously fails to describe satellite features. One might hope to remedy this by incorporating self-consistency into the perturbative expansion. However, self-consistent $GW$ does not necessarily yield a better description of the satellites than $G_0W_0$; this has been attributed to the neglect of vertex corrections in $GW$.
% TODO: cite self-consistent GW satellite studies

Coupled-cluster theory, a nonperturbative framework which has been established as the gold standard of ground-state molecular quantum chemistry when up to triple excitations are considered (CCSDT), offers a route to these vertex corrections. Excited-state properties, including spectral functions, may be obtained from its equation-of-motion extension (EOM-CC). However, the steep computational scaling of EOM-CCSDT limits its application to small supercells. Although one may slightly reduce the scaling by truncating the excitation hierarchy at doubles, it is unclear to what extent this incorporates the desired vertex corrections.

The cumulant expansion presents a viable alternative, as it incorporates a class of vertex corrections by resumming the perturbation series. The cumulant expansion was identified as the exact solution for a single fermion coupled to a bath of bosons. As a result, it qualitatively reproduces the core spectrum, among other features. The cumulant expansion is typically performed on top of a $G_0W_0$ calculation, an approach we refer to as $G_0W_0+C$. Consequently, $G_0W_0+C$ inherits the starting-point dependence of $G_0W_0$. Furthermore, studies of finite supercells of the electron gas have shown that $G_0W_0+C$ deviates sizeably from numerically exact results. Therefore, we see a need to explore cumulant approaches apart from $G_0W_0+C$. 
% TODO: cite finite-supercell benchmarks

% A remedy lies in the observation that core spectroscopy gives rise to an electronic polaron, in which the suddenly created hole couples to the bosonic degrees of freedom of the system. In this work, we exploit this observation by writing down a physically motivated Hamiltonian for the electronic polaron (albeit deriving a controlled approximation to the \emph{ab initio} problem) and solving for its Green's function. We are aided in this endeavor by the rapid development, in recent years, of techniques that provide near-exact solutions of polaron problems at reasonable computational cost.
% \subsubsection{TODO: Why do we study the electron gas at rs = 4; It is a good model for alkali metals, and this density is also interesting because of the outside correlation contribution. Such that we are no longer in the perturbative regime
% }

In Sec.~\ref{sec:theory}, after recounting $G_0W_0$+C, we introduce BWs2+C, a cumulant method based on a renormalized perturbation expansion in the bare Coulomb interaction $v$. We also study a self-consistent cumulant expansion.
% This is followed by our quasiboson Hamiltonian and the techniques we use to solve it, namely the self-consistent cumulant expansion and delocalized coupled-cluster theory. 
In Sec.~\ref{sec:results}, we use the aforementioned methods to compute spectral functions. Finally, in Sec.~\ref{sec:conclusions}, we discuss our findings and and propose future directions.

\section{Theory}
In what follows, $i,j,\ldots$ and $a,b,\ldots$ denote momentum states below and above the Fermi level, respectively, while $p,q,\ldots$ denote general momentum states. We work in the spin-restricted formalism. For a momentum $k$, all
cumulant methods determine the time-domain single-particle Green's function $G_k(t)$ as 
\label{sec:theory}
\begin{equation}
G_k(t) = G_k^{0}(t)\, e^{C_k(t)},
\label{eq:cumulant}
\end{equation}
with $G_k^{0}(t)$ the non-interacting Green's function and $C_k(t)$ the cumulant. 
\subsection{BWs2+C}
In the BWs2 partition,
\begin{equation}
G^{(0)}_k(\omega)
=
\big(\underbrace{ \omega - \varepsilon_k^{\rm HF}}_{\left(G^{\rm{HF}}_k(\omega)\right)^{-1}} -( \underbrace{\alpha_1 \mathbf{W}^O \oplus \alpha_2 \mathbf{W}^V)_{kk}}_{{W}_{kk}} \big)^{-1},
\end{equation}
 where $\varepsilon^{\rm HF}_k$ is the HF single-particle energy of momentum state $k$. Elements of the $\bm{W}$ regularizer are $W^O_{ii}=\frac{1}{2}\Sigma_{ii}^{(2),\rm 2p1h}[\bm{\tilde{\varepsilon}}](\omega = \tilde{\varepsilon}_i)$ and $W^V_{aa}=\frac{1}{2}\Sigma_{aa}^{(2),\rm{2h1p}}[\bm{\tilde{\varepsilon}}](\omega = \tilde{\varepsilon}_a)$. $\Sigma_{pp}^{(2),\rm ch}[\mathbf{\tilde{\varepsilon}}](\omega)$ is the diagonal element of a given channel (ch=2h1p or ch=2p1h) of the second-order Born self-energy $\Sigma^{(2)}$ \emph{evaluated at} the BWs2 single-particle energies $\bm{\tilde{\varepsilon}}$ (whose $p$th element is $\tilde{\varepsilon}_p$). 

The interpretation is that in order to determine $G^{(0)}_k$ for BWs2+C, a self-consistent optimization must first be performed, that requires repeated evaluation of the second-order self-energy, which scales as $O(N^5)$. 
% This has the function of incorporating dynamical screening that the standard Moller-Plesset partitioning lacks. 

% This is reminiscent of, for example, EOM-CCSD, which, in the EOM calculation, formally exhibits fifth-order scaling, but the ground-state CCSD calculation remains at sixth-order scaling. 
The second Born expansion of the Dyson equation takes the form
\begin{align}
G^{-1}_k(\omega)
&\approx \left( G^{\rm{HF}}_k(\omega)\right)^{-1} - {\Sigma}^{(2)}_{kk}[\bm{\varepsilon}^{\rm HF}](\omega). \label{eq:dyson_born}
\end{align}
The BWs2 partition, chooses a different partition,\ where the Dyson equation takes the form
\begin{align}
G^{-1}_k(\omega)
% &= \omega - (\mathbf{F} + \mathbf{W})
%    - \big[ \bm{\Sigma}(\omega) - \mathbf{W} \big] \\
&\approx ( G^{(0)}_k(\omega) \big)^{-1}
   - {\Sigma}^{\mathrm{eff}}_{kk}(\omega)\label{eq:dyson_bws2},
\end{align}
where  ${\Sigma}^{\mathrm{eff}}_{kk}(\omega) \equiv {\Sigma}^{(2)}_{kk}[\bm{\tilde{\varepsilon}}](\omega) - {W}_{kk}$. We are motivated to see how the MBPT truncation of \ref{eq:dyson_bws2} fares for computing excited-states in the solid-state because the renormalization of single-particle energies in BWs2 is hypothesized to have the function of incorporating incorporates dynamical screening. Furthermore, it has been shown by Tew that BWs2 facilitates the accurate computation of quasiparticle energies for strongly correlated molecular systems. Note that the approximation in \ref{eq:dyson_bws2} remains exact up to second order in $v$.
Equating the cumulant and Dyson expansions to second order in the bare Coulomb interaction gives 
\begin{align}
C_{k}(t)
&=
i \int \frac{d\omega}{2\pi}
e^{-i(\omega - \tilde{\varepsilon}_k)t}
\frac{\Sigma^{\mathrm{eff}}_{kk}(\omega)}
     {(\omega - \tilde{\varepsilon}_k)^2} \\
&=
i \int \frac{d\omega'}{2\pi}
e^{-i\omega' t}
\frac{\Sigma^{\mathrm{eff}}_{kk}(\omega' + \tilde{\varepsilon}_k)}
     {(\omega')^2}. %\\
% &=
% \int d\omega\,
% \frac{\beta_{pp}(\omega)}{\omega^2}
% \left(
% e^{-i \omega t}
% + i \omega t
% - 1
% \right),
\end{align}

For momentum $k<k_F$, we have
\begin{align}
  \Sigma^{\mathrm{eff}}_{k}(\omega) 
&= \Sigma_{kk}^{(2)}\!\left[ \bm{\tilde{\varepsilon}} \right](\omega) - \frac{\alpha_1}{2} \Sigma_{kk}^{(2),\rm{2p1h}}\!\left[ \bm{\tilde{\varepsilon}} \right](\tilde{\varepsilon}_k)\\
\implies \Sigma^{\mathrm{eff}}_{k}(\omega + \tilde{\varepsilon}_k)  &=
  \Sigma_{kk}^{(2)}\!\left[ \bm{\tilde{\varepsilon}} \right](\omega + \tilde{\varepsilon}_k) - \frac{\alpha_1}{2} \Sigma_{kk}^{(2),\rm 2p1h}\!\left[ \bm{\tilde{\varepsilon}} \right](\tilde{\varepsilon}_k)\\
\implies C_{k}(t) &= \sum_{iab} \frac{\mathcal{K}_{ki}^{ab} }{\left(\tilde{\epsilon}_{k i}^{a b}\right)^2}\left(e^{-i \tilde{\epsilon}_{k i}^{a b} t} +i \left( 1-\alpha_1 /2\right) \tilde{\epsilon}_{k i}^{a b} t -1 \right) \nonumber\\&+\sum_{ija} \frac{\mathcal{K}_{ij}^{ka}}{\left(\tilde{\epsilon}_{k a}^{i j}\right)^2} \left(e^{-i \tilde{\epsilon}_{k a}^{i j} t} +i \tilde{\epsilon}_{k a}^{i j} t -1 \right),
% \begin{equation}
%     C_{pp}^{(2;\alpha_1)}(t) = 
%     \sum_{i a b}\frac{\mathcal{K}_{pi}^{ab} }{\left(\tilde{\epsilon}_{p i}^{a b}\right)^2}\left(e^{-i \tilde{\epsilon}_{p i}^{a b} t} +i \left( 1-\alpha_1 /2\right) \tilde{\epsilon}_{p i}^{a b} t -1 \right) +\sum_{i j a} \frac{\mathcal{K}_{ij}^{pa}}{\left(\tilde{\epsilon}_{p a}^{i j}\right)^2} \left(e^{-i \tilde{\epsilon}_{p a}^{i j} t} +i \tilde{\epsilon}_{p a}^{i j} t -1 \right)
% \end{equation}
\end{align}
where $\mathcal{K}_{pq}^{rs} = (pr|qs) \left[ 2(pr|qs)-(ps|qr) \right]$,
$\tilde{\epsilon}_{ki}^{ab} = \tilde{\varepsilon}_a + \tilde{\varepsilon}_b - \tilde{\varepsilon}_k - \tilde{\varepsilon}_i$, and $\tilde{\epsilon}_{ka}^{ij} = \tilde{\varepsilon}_i + \tilde{\varepsilon}_j - \tilde{\varepsilon}_k - \tilde{\varepsilon}_a$. By analogy for momentum $k>k_F$, we have
\begin{align}
  C_{k}(t) &= \sum_{iab} \frac{\mathcal{K}_{ki}^{ab} }{\left(\tilde{\epsilon}_{k i}^{a b}\right)^2}\left(e^{-i \tilde{\epsilon}_{k i}^{a b} t} +i \tilde{\epsilon}_{k i}^{a b} t -1 \right) \nonumber\\&+\sum_{ija} \frac{\mathcal{K}_{ij}^{ka}}{\left(\tilde{\epsilon}_{k a}^{i j}\right)^2} \left(e^{-i \tilde{\epsilon}_{k a}^{i j} t} +i \left( 1-\alpha_2 /2\right) \tilde{\epsilon}_{k a}^{i j} t -1 \right).
\end{align}
Note that this reduces to the second-order cumulant when $\alpha_1 = \alpha_2 = 0$.
\subsection{GW+C}
We compute spectral functions of the homogeneous electron gas at a Wigner--Seitz radius of $r_s = 4$. The electron gas is a bellwether system for excited-state methods, whose ultimate goal is to predict the spectral functions of real materials. This density is of practical interest because it coincides with the valence electron density of the simple alkali metals. It is also of theoretical interest because, at $r_s = 4$, the electron gas lies outside the perturbative regime, where there is a notable lack of existing  methods that can accurately compute excited states.

The theories investigated below make use of the screened Coulomb interaction, in which the screening is obtained from the direct random-phase approximation (RPA). This amounts to solving the Casida equation
\begin{equation}
\begin{pmatrix}
\mathbf{A} & \mathbf{B} \\
-\mathbf{B} & -\mathbf{A}
\end{pmatrix}
\begin{pmatrix}
\mathbf{X} \\
\mathbf{Y}
\end{pmatrix}
=
\begin{pmatrix}
\mathbf{X} \\
\mathbf{Y}
\end{pmatrix}
\mathbf{\Omega},
\label{eq:casida}
%     \left(\begin{array}{ll}
% \mathrm{A} & \mathrm{~B} \\
% \mathrm{~B} & \mathrm{~A}
% \end{array}\right)\left(\begin{array}{ll}
% \mathrm{X} & \mathrm{Y} \\
% \mathrm{Y} & \mathrm{X}
% \end{array}\right)=\left(\begin{array}{cc}
% \mathrm{X} & -\mathrm{Y} \\
% -\mathrm{Y} & \mathrm{X}
% \end{array}\right)\left(\begin{array}{ll}
% \Omega & 0 \\
% 0 & \Omega
% \end{array}\right)
\end{equation}
where
\begin{align}
    A_{ia,jb} &= \delta_{ij}\delta_{ab} \left(\epsilon_a^{\rm LDA}-\epsilon_i^{\rm LDA}\right) + v(\mathbf{q}),\\
    B_{ia,jb} &= v(\mathbf{q}).
% \left(\green{\epsilon_a^{\rm HF}}+\red{\frac{v_M}{2}}-\green{\epsilon_i^{\rm HF}}-\red{\frac{v_M}{2}}\right)+\langle\Phi_i^a|\green{\hat V_{\rm res}}|\Phi_j^b\rangle\\
%     &= \delta_{ij}\delta_{ab}\left(\green{\epsilon_a^{\rm HF}}-\green{\epsilon_i^{\rm HF}}\right)+\langle\Phi_i^a|\green{\hat V_{\rm res}}|\Phi_j^b\rangle\\
% B_{ia,jb} &= \langle\Phi_i^a|\green{\hat V_{\rm res}}|\Phi_j^b\rangle.
\end{align}
Here, $i$ and $j$ label bands, $a = i + \mathbf{q}$ and $b = j + \mathbf{q}$ label unoccupied orbitals, and $v(\mathbf{q}) = 4\pi/(V q^2)$ is the Coulomb interaction in a supercell of volume $V$.

Because momentum is conserved, the Casida matrix is block diagonal in the transferred momentum $\mathbf{q}$. Diagonalizing each block yields a set of bosonic modes $\nu$ (which include the plasmon) with frequencies $\Omega_\nu(\mathbf{q})$. The corresponding hole--boson coupling is
\begin{equation}
g_{\nu}(\mathbf{k},\mathbf{q}) = \sqrt{2}\sum_{ia} (\mathbf{k}\,,\mathbf{k}-\mathbf{q}\,|\,i\,a)\,(\mathbf{X}+\mathbf{Y})_{ia,\nu}
= \sqrt{2}\,v(\mathbf{q})\sum_{i}(\mathbf{X}+\mathbf{Y})_{i,\,i+\mathbf{q};\,\nu},
\label{eq:coupling}
\end{equation}
where the second equality follows because, for plane waves, every allowed two-electron integral equals $v(\mathbf{q})$; the coupling is therefore independent of $\mathbf{k}$. A naive diagonalization of the full Casida matrix scales as the sixth power of system size. By exploiting momentum conservation, we instead diagonalize a set of much smaller, independent blocks, one per momentum sector $\mathbf{q}$; the cubic symmetry of the electron gas further reduces the number of symmetry-inequivalent sectors that must be treated.

Note that in the electron gas, the LDA exchange--correlation potential is a constant, so the LDA eigenvalues are simply the free-electron kinetic energies shifted by $V^{\rm xc}$. The LDA starting point therefore contains no correlation in its dispersion.
% \subsubsection{TODO: Introduce our RPA theory. }

\subsection{One-shot {GW}+C}
\label{sec:gwc}
In the cumulant formalism, the Green's function is written as
\begin{equation}
G_k(t) = G_k^{0}(t)\, e^{C_k(t)},
\label{eq:cumulant}
\end{equation}
where $G_k^{0}(t)$ is a non-interacting reference Green's function built from $\epsilon_k^{\rm LDA}$, and $C_k(t) = \tilde{C}_k(t) - i\Lambda_k t$ is the cumulant. The static part of the cumulant is $\Lambda_k = \Sigma_k^{\rm x} - V_k^{\rm xc}$, where $\Sigma_k^{\rm x}$ is the exchange self-energy. The reference Green's function and the dynamic part $\tilde{C}_k(t)$ depend on whether one constructs the time-ordered or the retarded cumulant, as described below.

\subsubsection{Time-ordered cumulant}
In the time-ordered cumulant (TC), the dynamic part is
\begin{equation}
    \tilde{C}_k^{\rm T}(t) = \int d\omega\, \frac{\beta_k^{\pm}(\omega)}{\omega^2}\left(e^{-i \omega t}+i \omega t-1\right),
    \label{eq:tc}
\end{equation}
where
\begin{equation}
\beta_k^{\pm}(\omega) = \frac{1}{\pi}\left|\operatorname{Im} \Sigma_k^{\rm R}\!\left(\omega+\epsilon_k^{\rm LDA}\right)\right| \theta\!\left(\pm\left(\omega + \epsilon_k^{\rm LDA} - \epsilon_{\rm F}\right)\right).
\end{equation}
The non-interacting reference is
\begin{equation}
G_k^{{\rm T},0}(t) = \mp i\, \theta(\pm t)\, \theta\!\left(\pm\left(\epsilon_k^{\rm LDA}-\epsilon_{\rm F}\right)\right) e^{-i \epsilon_k^{\rm LDA} t}.
\end{equation}
Throughout, the upper (lower) sign refers to particles with $k>k_{\rm F}$ (holes with $k<k_{\rm F}$), $\epsilon_{\rm F} \equiv \epsilon_{k_{\rm F}}^{\rm LDA}$ is the Fermi energy, and $\theta$ is the Heaviside step function. Thus, for a hole, only the portion of the self-energy below the Fermi level contributes to the cumulant.

\subsubsection{Retarded cumulant}
In the retarded cumulant (RC), the dynamic part is
\begin{equation}
    \tilde{C}_k^{\rm R}(t) = \int d\omega\, \frac{\beta_k^{\rm R}(\omega)}{\omega^2} \left( e^{-i\omega t} + i\omega t - 1 \right),
    \label{eq:rc}
\end{equation}
where
\begin{equation}
\beta_k^{\rm R}(\omega) = \frac{1}{\pi}\left|\operatorname{Im} \Sigma_k^{\rm R}\!\left(\omega+\epsilon_k^{\rm LDA}\right)\right|,
\end{equation}
so that both the particle and hole portions of the self-energy contribute. The non-interacting reference always takes the form
\begin{equation}
G_k^{{\rm R},0}(t) = -i\, \theta(t)\, e^{-i\epsilon_k^{\rm LDA} t}.
\end{equation}
Kas, Rehr, and Reining showed that, in the thermodynamic limit, the retarded cumulant yields a more physical spectral function because it treats particles and holes on an equal footing~\cite{kas2014}. They also showed that, at the bottom of the band, the TC and RC are essentially equivalent.
\subsection{Quasiboson formalism}
The methods discussed below solve the Hamiltonian
\begin{align}
\hat{H}
&=
\hat{H}_{\rm el}
+
\hat{H}_{\rm pl}
+
\hat{H}_{\rm ep}
\end{align}
with $\hat{H}_{\rm el}=  \sum_{\bm{k}\in\mathcal{O}}
\epsilon_{\bm{k}}^{\rm HF}
a_{\bm{k}}^{\dagger}a_{\bm{k}}$, $\hat{H}_{\rm pl}= \sum_{\nu_{\bm{q}}\in\mathcal{M}}
\omega_{\nu_{\bm{q}}}
b_{\nu_{\bm{q}}}^{\dagger}
b_{\nu_{\bm{q}}}$, and $\hat{H}_{\rm ep}= \sum_{\bm{k},\bm{k}+\bm{q}\in\mathcal{O}}
\sum_{\nu_{\bm{q}}\in\mathcal{M}}
g_{\nu}^{\rm pl}(\bm{k},\bm{q})
a_{\bm{k}+\bm{q}}^{\dagger}
a_{\bm{k}}
\left(
b_{\nu_{-\bm{q}}}^{\dagger}
+
b_{\nu_{\bm{q}}}
\right)$, where $\mathcal{O}$ and $\mathcal{M}$ are the sets of occupied bands and plasmonic modes, respectively. As indicated, given a plasmonic mode $\nu$, the associated momentum transfer $\bm{q}$ is already determined. The Plasmonic  frequencies $\omega_{\nu_{\bm{q}}}$and electron-plasmon couplings $g_{\nu}^{\rm pl}(\bm{k},\bm{q})$ are The eigenvalues  and transition densities taken from a direct RPA calculation starting from n. LDA reference as . The seasoned GW practitioner may be wary of this formulation because it seemingly mixes Hartree-fock single-particle energies with LDA screening. However, this has been done with empirical success in the self-consistent GW community for decades, so we pursue the same approach in defining the quasiboson Hamiltonian. 

\subsubsection{Self-consistent cumulant expansion}
For many decades, self-consistency has been pursued has been pursued within the Dyson framework. There have only been a few studies of carrying out the self-consistency in the cumulant formalism, which have shown promise for computing spectral functions of the hosting model . This is the first.  application of the self-consistent cumulant expansion to A fermionic system in the electron gas. Note that this has been done before by Holmes, but this is not true self-consistency, as we show in the appendix.  Specifically, we compute the retarded Green's function in the time domain as In equation X, and then we solve for the cumulant as 
We are always solving
\begin{equation}
\frac{dC_k}{dt}=-i\Lambda_k - \sum_{i\in\mathcal{S}}\sum_\nu|g_{ki\nu}|^2\int_0^t\!d\tau\; e^{-i\Delta_{ki\nu}(t-\tau)}\; e^{-[C_k(t)-C_k(\tau)]+[C_i(t)-C_i(\tau)]},
\end{equation}
where $\Delta_{ki\nu}=\epsilon^{\rm LDA}_i-\epsilon^{\rm LDA}_k-{\Omega_\nu}$, $\Lambda_k$ is a constant, ${g_{ki\nu}}$ is the electron-plasmon coupling, and $\mathcal{S}$ is some set of states.
\subsubsection{Delocalized coupled-cluster theory}
The D2 ansatz is 
We compute the DCC ground state by considering the following processes of excitation out of the reference Hartree-fock reference. Then we apply the tangent-space formalism to compute the spectral functions. We use the nomenclature dCC-X-SY-(Z, A) To denote a method, which, For the ground state calculation is X order In the fermions with Y coupled electron-boson excitations. See this paper for more details. 
\section{Results}
Here, we study finite supercells of the electron gas at rs = 4, converging onto the thermodynamic limit. To begin, we consider the smallest possible supercell with 14 electrons in 19 orbitals (14e19o), for which, We can compare the methods of our study with  the numerically exact result with DMRG, which well coincides wtith CCSDT.\\
For 114e485o, 
\section{Conclusions}
% \section{\label{sec:level1}First-level heading:\protect\\ The line
% break was forced \lowercase{via} \textbackslash\textbackslash}

% This sample document demonstrates proper use of REV\TeX~4.2 (and
% \LaTeXe) in mansucripts prepared for submission to APS
% journals. Further information can be found in the REV\TeX~4.2
% documentation included in the distribution or available at
% \url{http://journals.aps.org/revtex/}.

% When commands are referred to in this example file, they are always
% shown with their required arguments, using normal \TeX{} format. In
% this format, \verb+#1+, \verb+#2+, etc. stand for required
% author-supplied arguments to commands. For example, in
% \verb+\section{#1}+ the \verb+#1+ stands for the title text of the
% author's section heading, and in \verb+\title{#1}+ the \verb+#1+
% stands for the title text of the paper.

% Line breaks in section headings at all levels can be introduced using
% \textbackslash\textbackslash. A blank input line tells \TeX\ that the
% paragraph has ended. Note that top-level section headings are
% automatically uppercased. If a specific letter or word should appear in
% lowercase instead, you must escape it using \verb+\lowercase{#1}+ as
% in the word ``via'' above.

% \subsection{\label{sec:level2}Second-level heading: Formatting}

% This file may be formatted in either the \texttt{preprint} or
% \texttt{reprint} style. \texttt{reprint} format mimics final journal output. 
% Either format may be used for submission purposes. \texttt{letter} sized paper should
% be used when submitting to APS journals.

% \subsubsection{Wide text (A level-3 head)}
% The \texttt{widetext} environment will make the text the width of the
% full page, as on page~\pageref{eq:wideeq}. (Note the use the
% \verb+\pageref{#1}+ command to refer to the page number.) 
% \paragraph{Note (Fourth-level head is run in)}
% The width-changing commands only take effect in two-column formatting. 
% There is no effect if text is in a single column.

% \subsection{\label{sec:citeref}Citations and References}
% A citation in text uses the command \verb+\cite{#1}+ or
% \verb+\onlinecite{#1}+ and refers to an entry in the bibliography. 
% An entry in the bibliography is a reference to another document.

% \subsubsection{Citations}
% Because REV\TeX\ uses the \verb+natbib+ package of Patrick Daly, 
% the entire repertoire of commands in that package are available for your document;
% see the \verb+natbib+ documentation for further details. Please note that
% REV\TeX\ requires version 8.31a or later of \verb+natbib+.

% \paragraph{Syntax}
% The argument of \verb+\cite+ may be a single \emph{key}, 
% or may consist of a comma-separated list of keys.
% The citation \emph{key} may contain 
% letters, numbers, the dash (-) character, or the period (.) character. 
% New with natbib 8.3 is an extension to the syntax that allows for 
% a star (*) form and two optional arguments on the citation key itself.
% The syntax of the \verb+\cite+ command is thus (informally stated)
% \begin{quotation}\flushleft\leftskip1em
% \verb+\cite+ \verb+{+ \emph{key} \verb+}+, or\\
% \verb+\cite+ \verb+{+ \emph{optarg+key} \verb+}+, or\\
% \verb+\cite+ \verb+{+ \emph{optarg+key} \verb+,+ \emph{optarg+key}\ldots \verb+}+,
% \end{quotation}\noindent
% where \emph{optarg+key} signifies 
% \begin{quotation}\flushleft\leftskip1em
% \emph{key}, or\\
% \texttt{*}\emph{key}, or\\
% \texttt{[}\emph{pre}\texttt{]}\emph{key}, or\\
% \texttt{[}\emph{pre}\texttt{]}\texttt{[}\emph{post}\texttt{]}\emph{key}, or even\\
% \texttt{*}\texttt{[}\emph{pre}\texttt{]}\texttt{[}\emph{post}\texttt{]}\emph{key}.
% \end{quotation}\noindent
% where \emph{pre} and \emph{post} is whatever text you wish to place 
% at the beginning and end, respectively, of the bibliographic reference
% (see Ref.~[\onlinecite{witten2001}] and the two under Ref.~[\onlinecite{feyn54}]).
% (Keep in mind that no automatic space or punctuation is applied.)
% It is highly recommended that you put the entire \emph{pre} or \emph{post} portion 
% within its own set of braces, for example: 
% \verb+\cite+ \verb+{+ \texttt{[} \verb+{+\emph{text}\verb+}+\texttt{]}\emph{key}\verb+}+.
% The extra set of braces will keep \LaTeX\ out of trouble if your \emph{text} contains the comma (,) character.

% The star (*) modifier to the \emph{key} signifies that the reference is to be 
% merged with the previous reference into a single bibliographic entry, 
% a common idiom in APS and AIP articles (see below, Ref.~[\onlinecite{epr}]). 
% When references are merged in this way, they are separated by a semicolon instead of 
% the period (full stop) that would otherwise appear.

% \paragraph{Eliding repeated information}
% When a reference is merged, some of its fields may be elided: for example, 
% when the author matches that of the previous reference, it is omitted. 
% If both author and journal match, both are omitted.
% If the journal matches, but the author does not, the journal is replaced by \emph{ibid.},
% as exemplified by Ref.~[\onlinecite{epr}]. 
% These rules embody common editorial practice in APS and AIP journals and will only
% be in effect if the markup features of the APS and AIP Bib\TeX\ styles is employed.

% \paragraph{The options of the cite command itself}
% Please note that optional arguments to the \emph{key} change the reference in the bibliography, 
% not the citation in the body of the document. 
% For the latter, use the optional arguments of the \verb+\cite+ command itself:
% \verb+\cite+ \texttt{*}\allowbreak
% \texttt{[}\emph{pre-cite}\texttt{]}\allowbreak
% \texttt{[}\emph{post-cite}\texttt{]}\allowbreak
% \verb+{+\emph{key-list}\verb+}+.

% \subsubsection{Example citations}
% By default, citations are numerical\cite{Beutler1994}.
% Author-year citations are used when the journal is RMP. 
% To give a textual citation, use \verb+\onlinecite{#1}+: 
% Refs.~\onlinecite{[][{, and references therein}]witten2001,Bire82}. 
% By default, the \texttt{natbib} package automatically sorts your citations into numerical order and ``compresses'' runs of three or more consecutive numerical citations.
% REV\TeX\ provides the ability to automatically change the punctuation when switching between journal styles that provide citations in square brackets and those that use a superscript style instead. This is done through the \texttt{citeautoscript} option. For instance, the journal style \texttt{prb} automatically invokes this option because \textit{Physical 
% Review B} uses superscript-style citations. The effect is to move the punctuation, which normally comes after a citation in square brackets, to its proper position before the superscript. 
% To illustrate, we cite several together 
% \cite{[See the explanation of time travel in ]feyn54,*[The classical relativistic treatment of ][ is a relative classic]epr,witten2001,Berman1983,Davies1998,Bire82}, 
% and once again in different order (Refs.~\cite{epr,feyn54,Bire82,Berman1983,witten2001,Davies1998}). 
% Note that the citations were both compressed and sorted. Futhermore, running this sample file under the \texttt{prb} option will move the punctuation to the correct place.

% When the \verb+prb+ class option is used, the \verb+\cite{#1}+ command
% displays the reference's number as a superscript rather than in
% square brackets. Note that the location of the \verb+\cite{#1}+
% command should be adjusted for the reference style: the superscript
% references in \verb+prb+ style must appear after punctuation;
% otherwise the reference must appear before any punctuation. This
% sample was written for the regular (non-\texttt{prb}) citation style.
% The command \verb+\onlinecite{#1}+ in the \texttt{prb} style also
% displays the reference on the baseline.

% \subsubsection{References}
% A reference in the bibliography is specified by a \verb+\bibitem{#1}+ command
% with the same argument as the \verb+\cite{#1}+ command.
% \verb+\bibitem{#1}+ commands may be crafted by hand or, preferably,
% generated by Bib\TeX. 
% REV\TeX~4.2 includes Bib\TeX\ style files
% \verb+apsrev4-2.bst+, \verb+apsrmp4-2.bst+ appropriate for
% \textit{Physical Review} and \textit{Reviews of Modern Physics},
% respectively.

% \subsubsection{Example references}
% This sample file employs the \verb+\bibliography+ command, 
% which formats the \texttt{\jobname .bbl} file
% and specifies which bibliographic databases are to be used by Bib\TeX\ 
% (one of these should be by arXiv convention \texttt{\jobname .bib}).
% Running Bib\TeX\ (via \texttt{bibtex \jobname}) 
% after the first pass of \LaTeX\ produces the file
% \texttt{\jobname .bbl} which contains the automatically formatted
% \verb+\bibitem+ commands (including extra markup information via
% \verb+\bibinfo+ and \verb+\bibfield+ commands). 
% If not using Bib\TeX, you will have to create the \verb+thebibiliography+ environment 
% and its \verb+\bibitem+ commands by hand.

% Numerous examples of the use of the APS bibliographic entry types appear in the bibliography of this sample document.
% You can refer to the \texttt{\jobname .bib} file, 
% and compare its information to the formatted bibliography itself.

% \subsection{Footnotes}%
% Footnotes, produced using the \verb+\footnote{#1}+ command, 
% usually integrated into the bibliography alongside the other entries.
% Numerical citation styles do this%
% \footnote{Automatically placing footnotes into the bibliography requires using BibTeX to compile the bibliography.};
% author-year citation styles place the footnote at the bottom of the text column.
% Note: due to the method used to place footnotes in the bibliography, 
% \emph{you must re-run Bib\TeX\ every time you change any of your document's footnotes}. 

% \section{Math and Equations}
% Inline math may be typeset using the \verb+$+ delimiters. Bold math
% symbols may be achieved using the \verb+bm+ package and the
% \verb+\bm{#1}+ command it supplies. For instance, a bold $\alpha$ can
% be typeset as \verb+$\bm{\alpha}$+ giving $\bm{\alpha}$. Fraktur and
% Blackboard (or open face or double struck) characters should be
% typeset using the \verb+\mathfrak{#1}+ and \verb+\mathbb{#1}+ commands
% respectively. Both are supplied by the \texttt{amssymb} package. For
% example, \verb+$\mathbb{R}$+ gives $\mathbb{R}$ and
% \verb+$\mathfrak{G}$+ gives $\mathfrak{G}$

% In \LaTeX\ there are many different ways to display equations, and a
% few preferred ways are noted below. Displayed math will center by
% default. Use the class option \verb+fleqn+ to flush equations left.

% Below we have numbered single-line equations; this is the most common
% type of equation in \textit{Physical Review}:
% \begin{eqnarray}
% \chi_+(p)\alt{\bf [}2|{\bf p}|(|{\bf p}|+p_z){\bf ]}^{-1/2}
% \left(
% \begin{array}{c}
% |{\bf p}|+p_z\\
% px+ip_y
% \end{array}\right)\;,
% \\
% \left\{%
%  \openone234567890abc123\alpha\beta\gamma\delta1234556\alpha\beta
%  \frac{1\sum^{a}_{b}}{A^2}%
% \right\}%
% \label{eq:one}.
% \end{eqnarray}
% Note the open one in Eq.~(\ref{eq:one}).

% Not all numbered equations will fit within a narrow column this
% way. The equation number will move down automatically if it cannot fit
% on the same line with a one-line equation:
% \begin{equation}
% \left\{
%  ab12345678abc123456abcdef\alpha\beta\gamma\delta1234556\alpha\beta
%  \frac{1\sum^{a}_{b}}{A^2}%
% \right\}.
% \end{equation}

% When the \verb+\label{#1}+ command is used [cf. input for
% Eq.~(\ref{eq:one})], the equation can be referred to in text without
% knowing the equation number that \TeX\ will assign to it. Just
% use \verb+\ref{#1}+, where \verb+#1+ is the same name that used in
% the \verb+\label{#1}+ command.

% Unnumbered single-line equations can be typeset
% using the \verb+\[+, \verb+\]+ format:
% \[g^+g^+ \rightarrow g^+g^+g^+g^+ \dots ~,~~q^+q^+\rightarrow
% q^+g^+g^+ \dots ~. \]


% \subsection{Multiline equations}

% Multiline equations are obtained by using the \verb+eqnarray+
% environment.  Use the \verb+\nonumber+ command at the end of each line
% to avoid assigning a number:
% \begin{eqnarray}
% {\cal M}=&&ig_Z^2(4E_1E_2)^{1/2}(l_i^2)^{-1}
% \delta_{\sigma_1,-\sigma_2}
% (g_{\sigma_2}^e)^2\chi_{-\sigma_2}(p_2)\nonumber\\
% &&\times
% [\epsilon_jl_i\epsilon_i]_{\sigma_1}\chi_{\sigma_1}(p_1),
% \end{eqnarray}
% \begin{eqnarray}
% \sum \vert M^{\text{viol}}_g \vert ^2&=&g^{2n-4}_S(Q^2)~N^{n-2}
%         (N^2-1)\nonumber \\
%  & &\times \left( \sum_{i<j}\right)
%   \sum_{\text{perm}}
%  \frac{1}{S_{12}}
%  \frac{1}{S_{12}}
%  \sum_\tau c^f_\tau~.
% \end{eqnarray}
% \textbf{Note:} Do not use \verb+\label{#1}+ on a line of a multiline
% equation if \verb+\nonumber+ is also used on that line. Incorrect
% cross-referencing will result. Notice the use \verb+\text{#1}+ for
% using a Roman font within a math environment.

% To set a multiline equation without \emph{any} equation
% numbers, use the \verb+\begin{eqnarray*}+,
% \verb+\end{eqnarray*}+ format:
% \begin{eqnarray*}
% \sum \vert M^{\text{viol}}_g \vert ^2&=&g^{2n-4}_S(Q^2)~N^{n-2}
%         (N^2-1)\\
%  & &\times \left( \sum_{i<j}\right)
%  \left(
%   \sum_{\text{perm}}\frac{1}{S_{12}S_{23}S_{n1}}
%  \right)
%  \frac{1}{S_{12}}~.
% \end{eqnarray*}

% To obtain numbers not normally produced by the automatic numbering,
% use the \verb+\tag{#1}+ command, where \verb+#1+ is the desired
% equation number. For example, to get an equation number of
% (\ref{eq:mynum}),
% \begin{equation}
% g^+g^+ \rightarrow g^+g^+g^+g^+ \dots ~,~~q^+q^+\rightarrow
% q^+g^+g^+ \dots ~. \tag{2.6$'$}\label{eq:mynum}
% \end{equation}

% \paragraph{A few notes on \texttt{tag}s} 
% \verb+\tag{#1}+ requires the \texttt{amsmath} package. 
% Place the \verb+\tag{#1}+ command before the \verb+\label{#1}+, if any. 
% The numbering produced by \verb+\tag{#1}+ \textit{does not affect} 
% the automatic numbering in REV\TeX; 
% therefore, the number must be known ahead of time, 
% and it must be manually adjusted if other equations are added. 
% \verb+\tag{#1}+ works with both single-line and multiline equations. 
% \verb+\tag{#1}+ should only be used in exceptional cases---%
% do not use it to number many equations in your paper. 
% Please note that this feature of the \texttt{amsmath} package
% is \emph{not} compatible with the \texttt{hyperref} (6.77u) package.

% Enclosing display math within
% \verb+\begin{subequations}+ and \verb+\end{subequations}+ will produce
% a set of equations that are labeled with letters, as shown in
% Eqs.~(\ref{subeq:1}) and (\ref{subeq:2}) below.
% You may include any number of single-line and multiline equations,
% although it is probably not a good idea to follow one display math
% directly after another.
% \begin{subequations}
% \label{eq:whole}
% \begin{eqnarray}
% {\cal M}=&&ig_Z^2(4E_1E_2)^{1/2}(l_i^2)^{-1}
% (g_{\sigma_2}^e)^2\chi_{-\sigma_2}(p_2)\nonumber\\
% &&\times
% [\epsilon_i]_{\sigma_1}\chi_{\sigma_1}(p_1).\label{subeq:2}
% \end{eqnarray}
% \begin{equation}
% \left\{
%  abc123456abcdef\alpha\beta\gamma\delta1234556\alpha\beta
%  \frac{1\sum^{a}_{b}}{A^2}
% \right\},\label{subeq:1}
% \end{equation}
% \end{subequations}
% Giving a \verb+\label{#1}+ command directly after the \verb+\begin{subequations}+, 
% allows you to reference all the equations in the \texttt{subequations} environment. 
% For example, the equations in the preceding subequations environment were
% Eqs.~(\ref{eq:whole}).

% \subsubsection{Wide equations}
% The equation that follows is set in a wide format, i.e., it spans the full page. 
% The wide format is reserved for long equations
% that cannot easily be set in a single column:
% \begin{widetext}
% \begin{equation}
% {\cal R}^{(\text{d})}=
%  g_{\sigma_2}^e
%  \left(
%    \frac{[\Gamma^Z(3,21)]_{\sigma_1}}{Q_{12}^2-M_W^2}
%   +\frac{[\Gamma^Z(13,2)]_{\sigma_1}}{Q_{13}^2-M_W^2}
%  \right)
%  + x_WQ_e
%  \left(
%    \frac{[\Gamma^\gamma(3,21)]_{\sigma_1}}{Q_{12}^2-M_W^2}
%   +\frac{[\Gamma^\gamma(13,2)]_{\sigma_1}}{Q_{13}^2-M_W^2}
%  \right)\;. 
%  \label{eq:wideeq}
% \end{equation}
% \end{widetext}
% This is typed to show how the output appears in wide format.
% (Incidentally, since there is no blank line between the \texttt{equation} environment above 
% and the start of this paragraph, this paragraph is not indented.)

% \section{Cross-referencing}
% REV\TeX{} will automatically number such things as
% sections, footnotes, equations, figure captions, and table captions. 
% In order to reference them in text, use the
% \verb+\label{#1}+ and \verb+\ref{#1}+ commands. 
% To reference a particular page, use the \verb+\pageref{#1}+ command.

% The \verb+\label{#1}+ should appear 
% within the section heading, 
% within the footnote text, 
% within the equation, or 
% within the table or figure caption. 
% The \verb+\ref{#1}+ command
% is used in text at the point where the reference is to be displayed.  
% Some examples: Section~\ref{sec:level1} on page~\pageref{sec:level1},
% Table~\ref{tab:table1},%
% \begin{table}[b]%The best place to locate the table environment is directly after its first reference in text
% \caption{\label{tab:table1}%
% A table that fits into a single column of a two-column layout. 
% Note that REV\TeX~4 adjusts the intercolumn spacing so that the table fills the
% entire width of the column. Table captions are numbered
% automatically. 
% This table illustrates left-, center-, decimal- and right-aligned columns,
% along with the use of the \texttt{ruledtabular} environment which sets the 
% Scotch (double) rules above and below the alignment, per APS style.
% }
% \begin{ruledtabular}
% \begin{tabular}{lcdr}
% \textrm{Left\footnote{Note a.}}&
% \textrm{Centered\footnote{Note b.}}&
% \multicolumn{1}{c}{\textrm{Decimal}}&
% \textrm{Right}\\
% \colrule
% 1 & 2 & 3.001 & 4\\
% 10 & 20 & 30 & 40\\
% 100 & 200 & 300.0 & 400\\
% \end{tabular}
% \end{ruledtabular}
% \end{table}
% and Fig.~\ref{fig:epsart}.%
% \begin{figure}[b]
% \includegraphics{fig_1}% Here is how to import EPS art
% \caption{\label{fig:epsart} A figure caption. The figure captions are
% automatically numbered.}
% \end{figure}

% \section{Floats: Figures, Tables, Videos, etc.}
% Figures and tables are usually allowed to ``float'', which means that their
% placement is determined by \LaTeX, while the document is being typeset. 

% Use the \texttt{figure} environment for a figure, the \texttt{table} environment for a table.
% In each case, use the \verb+\caption+ command within to give the text of the
% figure or table caption along with the \verb+\label+ command to provide
% a key for referring to this figure or table.
% The typical content of a figure is an image of some kind; 
% that of a table is an alignment.%
% \begin{figure*}
% \includegraphics{fig_2}% Here is how to import EPS art
% \caption{\label{fig:wide}Use the figure* environment to get a wide
% figure that spans the page in \texttt{twocolumn} formatting.}
% \end{figure*}
% \begin{table*}
% \caption{\label{tab:table3}This is a wide table that spans the full page
% width in a two-column layout. It is formatted using the
% \texttt{table*} environment. It also demonstates the use of
% \textbackslash\texttt{multicolumn} in rows with entries that span
% more than one column.}
% \begin{ruledtabular}
% \begin{tabular}{ccccc}
%  &\multicolumn{2}{c}{$D_{4h}^1$}&\multicolumn{2}{c}{$D_{4h}^5$}\\
%  Ion&1st alternative&2nd alternative&lst alternative
% &2nd alternative\\ \hline
%  K&$(2e)+(2f)$&$(4i)$ &$(2c)+(2d)$&$(4f)$ \\
%  Mn&$(2g)$\footnote{The $z$ parameter of these positions is $z\sim\frac{1}{4}$.}
%  &$(a)+(b)+(c)+(d)$&$(4e)$&$(2a)+(2b)$\\
%  Cl&$(a)+(b)+(c)+(d)$&$(2g)$\footnotemark[1]
%  &$(4e)^{\text{a}}$\\
%  He&$(8r)^{\text{a}}$&$(4j)^{\text{a}}$&$(4g)^{\text{a}}$\\
%  Ag& &$(4k)^{\text{a}}$& &$(4h)^{\text{a}}$\\
% \end{tabular}
% \end{ruledtabular}
% \end{table*}

% Insert an image using either the \texttt{graphics} or
% \texttt{graphix} packages, which define the \verb+\includegraphics{#1}+ command.
% (The two packages differ in respect of the optional arguments 
% used to specify the orientation, scaling, and translation of the image.) 
% To create an alignment, use the \texttt{tabular} environment. 

% The best place to locate the \texttt{figure} or \texttt{table} environment
% is immediately following its first reference in text; this sample document
% illustrates this practice for Fig.~\ref{fig:epsart}, which
% shows a figure that is small enough to fit in a single column. 

% In exceptional cases, you will need to move the float earlier in the document, as was done
% with Table~\ref{tab:table3}: \LaTeX's float placement algorithms need to know
% about a full-page-width float earlier. 

% Fig.~\ref{fig:wide}
% has content that is too wide for a single column,
% so the \texttt{figure*} environment has been used.%
% \begin{table}[b]
% \caption{\label{tab:table4}%
% Numbers in columns Three--Five are aligned with the ``d'' column specifier 
% (requires the \texttt{dcolumn} package). 
% Non-numeric entries (those entries without a ``.'') in a ``d'' column are aligned on the decimal point. 
% Use the ``D'' specifier for more complex layouts. }
% \begin{ruledtabular}
% \begin{tabular}{ccddd}
% One&Two&
% \multicolumn{1}{c}{\textrm{Three}}&
% \multicolumn{1}{c}{\textrm{Four}}&
% \multicolumn{1}{c}{\textrm{Five}}\\
% %\mbox{Three}&\mbox{Four}&\mbox{Five}\\
% \hline
% one&two&\mbox{three}&\mbox{four}&\mbox{five}\\
% He&2& 2.77234 & 45672. & 0.69 \\
% C\footnote{Some tables require footnotes.}
%   &C\footnote{Some tables need more than one footnote.}
%   & 12537.64 & 37.66345 & 86.37 \\
% \end{tabular}
% \end{ruledtabular}
% \end{table}

% The content of a table is typically a \texttt{tabular} environment, 
% giving rows of type in aligned columns. 
% Column entries separated by \verb+&+'s, and 
% each row ends with \textbackslash\textbackslash. 
% The required argument for the \texttt{tabular} environment
% specifies how data are aligned in the columns. 
% For instance, entries may be centered, left-justified, right-justified, aligned on a decimal
% point. 
% Extra column-spacing may be be specified as well, 
% although REV\TeX~4 sets this spacing so that the columns fill the width of the
% table. Horizontal rules are typeset using the \verb+\hline+
% command. The doubled (or Scotch) rules that appear at the top and
% bottom of a table can be achieved enclosing the \texttt{tabular}
% environment within a \texttt{ruledtabular} environment. Rows whose
% columns span multiple columns can be typeset using the
% \verb+\multicolumn{#1}{#2}{#3}+ command (for example, see the first
% row of Table~\ref{tab:table3}).%

% Tables~\ref{tab:table1}, \ref{tab:table3}, \ref{tab:table4}, and \ref{tab:table2}%
% \begin{table}[b]
% \caption{\label{tab:table2}
% A table with numerous columns that still fits into a single column. 
% Here, several entries share the same footnote. 
% Inspect the \LaTeX\ input for this table to see exactly how it is done.}
% \begin{ruledtabular}
% \begin{tabular}{cccccccc}
%  &$r_c$ (\AA)&$r_0$ (\AA)&$\kappa r_0$&
%  &$r_c$ (\AA) &$r_0$ (\AA)&$\kappa r_0$\\
% \hline
% Cu& 0.800 & 14.10 & 2.550 &Sn\footnotemark[1]
% & 0.680 & 1.870 & 3.700 \\
% Ag& 0.990 & 15.90 & 2.710 &Pb\footnotemark[2]
% & 0.450 & 1.930 & 3.760 \\
% Au& 1.150 & 15.90 & 2.710 &Ca\footnotemark[3]
% & 0.750 & 2.170 & 3.560 \\
% Mg& 0.490 & 17.60 & 3.200 &Sr\footnotemark[4]
% & 0.900 & 2.370 & 3.720 \\
% Zn& 0.300 & 15.20 & 2.970 &Li\footnotemark[2]
% & 0.380 & 1.730 & 2.830 \\
% Cd& 0.530 & 17.10 & 3.160 &Na\footnotemark[5]
% & 0.760 & 2.110 & 3.120 \\
% Hg& 0.550 & 17.80 & 3.220 &K\footnotemark[5]
% &  1.120 & 2.620 & 3.480 \\
% Al& 0.230 & 15.80 & 3.240 &Rb\footnotemark[3]
% & 1.330 & 2.800 & 3.590 \\
% Ga& 0.310 & 16.70 & 3.330 &Cs\footnotemark[4]
% & 1.420 & 3.030 & 3.740 \\
% In& 0.460 & 18.40 & 3.500 &Ba\footnotemark[5]
% & 0.960 & 2.460 & 3.780 \\
% Tl& 0.480 & 18.90 & 3.550 & & & & \\
% \end{tabular}
% \end{ruledtabular}
% \footnotetext[1]{Here's the first, from Ref.~\onlinecite{feyn54}.}
% \footnotetext[2]{Here's the second.}
% \footnotetext[3]{Here's the third.}
% \footnotetext[4]{Here's the fourth.}
% \footnotetext[5]{And etc.}
% \end{table}
% show various effects.
% A table that fits in a single column employs the \texttt{table}
% environment. 
% Table~\ref{tab:table3} is a wide table, set with the \texttt{table*} environment. 
% Long tables may need to break across pages. 
% The most straightforward way to accomplish this is to specify
% the \verb+[H]+ float placement on the \texttt{table} or
% \texttt{table*} environment. 
% However, the \LaTeXe\ package \texttt{longtable} allows headers and footers to be specified for each page of the table. 
% A simple example of the use of \texttt{longtable} can be found
% in the file \texttt{summary.tex} that is included with the REV\TeX~4
% distribution.

% There are two methods for setting footnotes within a table (these
% footnotes will be displayed directly below the table rather than at
% the bottom of the page or in the bibliography). The easiest
% and preferred method is just to use the \verb+\footnote{#1}+
% command. This will automatically enumerate the footnotes with
% lowercase roman letters. However, it is sometimes necessary to have
% multiple entries in the table share the same footnote. In this case,
% there is no choice but to manually create the footnotes using
% \verb+\footnotemark[#1]+ and \verb+\footnotetext[#1]{#2}+.
% \texttt{\#1} is a numeric value. Each time the same value for
% \texttt{\#1} is used, the same mark is produced in the table. The
% \verb+\footnotetext[#1]{#2}+ commands are placed after the \texttt{tabular}
% environment. Examine the \LaTeX\ source and output for
% Tables~\ref{tab:table1} and \ref{tab:table2}
% for examples.

% Video~\ref{vid:PRSTPER.4.010101} 
% illustrates several features new with REV\TeX4.2,
% starting with the \texttt{video} environment, which is in the same category with
% \texttt{figure} and \texttt{table}.%
% \begin{video}
% \href{http://prst-per.aps.org/multimedia/PRSTPER/v4/i1/e010101/e010101_vid1a.mpg}{\includegraphics{vid_1a}}%
%  \quad
% \href{http://prst-per.aps.org/multimedia/PRSTPER/v4/i1/e010101/e010101_vid1b.mpg}{\includegraphics{vid_1b}}
%  \setfloatlink{http://link.aps.org/multimedia/PRSTPER/v4/i1/e010101}%
%  \caption{\label{vid:PRSTPER.4.010101}%
%   Students explain their initial idea about Newton's third law to a teaching assistant. 
%   Clip (a): same force.
%   Clip (b): move backwards.
%  }%
% \end{video}
% The \verb+\setfloatlink+ command causes the title of the video to be a hyperlink to the
% indicated URL; it may be used with any environment that takes the \verb+\caption+
% command.
% The \verb+\href+ command has the same significance as it does in the context of
% the \texttt{hyperref} package: the second argument is a piece of text to be 
% typeset in your document; the first is its hyperlink, a URL.

% \textit{Physical Review} style requires that the initial citation of
% figures or tables be in numerical order in text, so don't cite
% Fig.~\ref{fig:wide} until Fig.~\ref{fig:epsart} has been cited.

% \begin{acknowledgments}
% We wish to acknowledge the support of the author community in using
% REV\TeX{}, offering suggestions and encouragement, testing new versions,
% \dots.
% \end{acknowledgments}

% \appendix

% \section{Appendixes}

% To start the appendixes, use the \verb+\appendix+ command.
% This signals that all following section commands refer to appendixes
% instead of regular sections. Therefore, the \verb+\appendix+ command
% should be used only once---to setup the section commands to act as
% appendixes. Thereafter normal section commands are used. The heading
% for a section can be left empty. For example,
% \begin{verbatim}
% \appendix
% \section{}
% \end{verbatim}
% will produce an appendix heading that says ``APPENDIX A'' and
% \begin{verbatim}
% \appendix
% \section{Background}
% \end{verbatim}
% will produce an appendix heading that says ``APPENDIX A: BACKGROUND''
% (note that the colon is set automatically).

% If there is only one appendix, then the letter ``A'' should not
% appear. This is suppressed by using the star version of the appendix
% command (\verb+\appendix*+ in the place of \verb+\appendix+).

% \section{A little more on appendixes}

% Observe that this appendix was started by using
% \begin{verbatim}
% \section{A little more on appendixes}
% \end{verbatim}

% Note the equation number in an appendix:
% \begin{equation}
% E=mc^2.
% \end{equation}

% \subsection{\label{app:subsec}A subsection in an appendix}

% You can use a subsection or subsubsection in an appendix. Note the
% numbering: we are now in Appendix~\ref{app:subsec}.

% Note the equation numbers in this appendix, produced with the
% subequations environment:
% \begin{subequations}
% \begin{eqnarray}
% E&=&mc, \label{appa}
% \\
% E&=&mc^2, \label{appb}
% \\
% E&\agt& mc^3. \label{appc}
% \end{eqnarray}
% \end{subequations}
% They turn out to be Eqs.~(\ref{appa}), (\ref{appb}), and (\ref{appc}).

% % The \nocite command causes all entries in a bibliography to be printed out
% % whether or not they are actually referenced in the text. This is appropriate
% % for the sample file to show the different styles of references, but authors
% % most likely will not want to use it.
% \nocite{*}

\bibliography{citations}

\end{document}
%
% ****** End of file apssamp.tex ******
